\documentclass[10pt]{article}
\usepackage[margin=1in]{geometry}
\usepackage{amsmath,amssymb,amsthm,booktabs,hyperref}
\newtheorem{theorem}{Theorem}
\newtheorem{lemma}{Lemma}
\newcommand{\coeff}[1]{[#1]}
\title{A nonunimodal line in a colored-partition coefficient array}
\author{Principia Math in Progress}
\date{12 September 2026}
\begin{document}
\maketitle

\paragraph{Task 1.}
MathDB 345259.
Keith's multicolor extension of a conjecture of Dousse and Kim asserts
that, after fixing the partition weight and setting one color variable
equal to \(1\), every lattice-line cut of the remaining coefficient array
is unimodal \cite[Conjecture 1]{Keith}. A color is assigned to each
\emph{distinct part size}; all copies of that size have the same color.
We disprove this extension with four colors.

\begin{theorem}
Take partitions of \(120\) with at most \(15\) parts, each at most \(15\).
Let \(H(x,y,z)\) be their four-color generating polynomial after setting
the fourth color variable equal to \(1\).
The coefficients at the consecutive lattice points
\[
 (11,1,1),\qquad(10,2,2),\qquad(9,3,3)
\]
on a line of primitive direction \((-1,1,1)\) are respectively
\[
 \boxed{95004,\qquad 90090,\qquad 100100.}
\]
In particular the line cut is not unimodal.
\end{theorem}

\begin{proof}
Write \(f_r\) for the number of uncolored partitions in this box and of
this weight with exactly \(r\) distinct part sizes. Assigning one of four
colors independently to each distinct size gives
\begin{equation}\label{H}
 H(x,y,z)=\sum_{r=0}^{15} f_r(1+x+y+z)^r.
\end{equation}
Only \(f_{13},f_{14},f_{15}\) can contribute to the three coefficients
in the theorem. We establish
\begin{equation}\label{f}
 f_{13}=504,\qquad f_{14}=0,\qquad f_{15}=1.
\end{equation}

The sum of the integers from \(1\) to \(15\) is \(120\). Thus there is
exactly one partition with \(15\) distinct sizes, namely
\(1+2+\cdots+15\).
For \(14\) distinct sizes, let \(a\) be the missing size.
The sum of one copy of each available size is \(120-a\).
At most one extra part is allowed. To recover weight \(120\), that part
would have to be \(a\), contrary to its being a missing size. Hence
\(f_{14}=0\).

For \(13\) distinct sizes, let \(1\leq a<b\leq15\) be the two missing
sizes. The base weight is \(120-a-b\); one or two additional parts
must have sum \(a+b\), and neither may be a missing size.
For \(3\leq s\leq29\), put
\[
 L_s=\#\{(a,b):1\leq a<b\leq15,\ a+b=s\}
 =\left\lfloor\frac{s-1}{2}\right\rfloor-\max(1,s-15)+1.
\]
If there is one extra part, it must be \(s=a+b\leq15\).
It is automatically different from \(a,b\), giving
\[
 \sum_{s=3}^{15}L_s
 =2(1+2+3+4+5+6)+7=49
\]
partitions of length \(14\).

For two extra parts \(c\leq d\), the number of pairs in
\(\{1,\ldots,15\}\) summing to \(s\) is
\(L_s+\mathbf 1_{2\mid s}\). A pair sharing a value with
\(\{a,b\}\) must be \((a,b)\) itself, so exactly one pair is forbidden.
Consequently the number of partitions of length \(15\) is
\[
 \sum_{s=3}^{29}L_s\bigl(L_s-1+\mathbf 1_{2\mid s}\bigr)=455.
\]
For completeness, reflection gives \(L_s=L_{32-s}\), \(L_{16}=7\),
and the sum over \(3\leq s\leq15\) is \(203\); thus the last
sum is \(2\cdot203+7^2=455\).
The missing pair and extra multiset recover each partition uniquely,
so no partition is counted twice. This proves \(f_{13}=49+455=504\).

By the multinomial theorem, for \(u+v+w\leq15\),
\begin{equation}\label{coeff}
 \coeff{x^uy^vz^w}H
 =\sum_{r=u+v+w}^{15}
   f_r\,\frac{r!}{u!\,v!\,w!\,(r-u-v-w)!}.
\end{equation}
Equations \eqref{f} and \eqref{coeff} now give
\begin{align*}
 \coeff{x^{11}yz}H
   &=504\,\frac{13!}{11!}
     +\frac{15!}{11!\,2!}
     =78624+16380=95004,\\
 \coeff{x^{10}y^2z^2}H
   &=\frac{15!}{10!\,2!\,2!}=90090,\\
 \coeff{x^9y^3z^3}H
   &=\frac{15!}{9!\,3!\,3!}=100100.
\end{align*}
The integer points of the relevant line have the form
\((12-j,j,j)\), with consecutive integers \(j\). The displayed points
correspond to \(j=1,2,3\). All are in the support.
Every contiguous subsequence of a unimodal sequence is unimodal, whereas
\(95004>90090<100100\). This proves the theorem.
\end{proof}

\paragraph{Scope.}
This disproves the universal multicolor statement. The original
two-color overpartition conjecture is a narrower question and is not
settled by this example. The argument uses the intended number of
remaining variables, \(k-1\); the source's dimensional typo plays
no role.

\paragraph{Verification.}
The accompanying standard-library Python verifier independently computes
the partition counts by a finite generating-function recurrence and
enumerates the missing-size constructions. All arithmetic is integral.
The proof above supplies the complete counting argument without relying
on a numerical search or on the verifier's output.

\paragraph{Task 2.}
Determine whether the multicolor conjecture holds with three colors,
so that the remaining array is two-dimensional.

\paragraph{Task 3.}
Determine the least partition weight and the smallest box admitting a
counterexample with four colors, specifying the ordering used to compare
boxes.

\paragraph{Task 4.}
Construct an infinite family of nonunimodal line cuts for four colors,
using partitions near the staircase weight in square boxes.

\paragraph{Task 5.}
Identify classes of line directions for which unimodality survives;
in particular, investigate coordinate lines and directions whose
nonzero coordinates all have the same sign.

\begin{thebibliography}{9}
\bibitem{Keith}
W.~J. Keith, \emph{Restricted \(k\)-color partitions, II},
International Journal of Number Theory \textbf{17} (2021), 591--601,
\href{https://doi.org/10.1142/S1793042120400151}
{doi:10.1142/S1793042120400151}.
Conjecture 1 and Theorem 3 are on page 2 of
\href{https://arxiv.org/abs/2001.08351}{arXiv:2001.08351v1};
the author's revised manuscript uploaded on 26 July 2021 also retains
that conjecture.
\end{thebibliography}
\end{document}
